\section{ĐÁNH GIÁ THỰC TRẠNG VÀ ĐỀ XUẤT GIẢI PHÁP CẢI THIỆN}

\subsection{Nhận xét điểm đạt (Ưu điểm) trong hoạt động quản lý và quản trị}
Qua quá trình phân tích thực trạng tại Công ty Cổ phần FPT, nhóm nghiên cứu nhận thấy doanh nghiệp đã xây dựng được nền móng quản lý và quản trị rất vững vàng với các điểm sáng nổi bật:

\begin{enumerate}
    \item \textbf{Tầm nhìn chiến lược dài hạn và đón đầu xu thế công nghệ:} Ban lãnh đạo FPT thể hiện năng lực hoạch định chiến lược xuất sắc khi liên tục chuyển mình qua các giai đoạn: từ phân phối máy tính sang gia công phần mềm, từ chuyển đổi số sang công nghệ bán dẫn và Trí tuệ Nhân tạo (chiến lược DC5).
    \item \textbf{Cơ cấu quản trị tiệm cận chuẩn mực quốc tế:} Việc tách bạch chức danh Chủ tịch HĐQT và Tổng Giám đốc từ năm 2013, duy trì tỷ lệ thành viên HĐQT độc lập đạt hơn 42\% với các chuyên gia quốc tế hàng đầu (từ Nhật Bản, Pháp, Ấn Độ) giúp các quyết sách chiến lược đạt tính khách quan và minh bạch cao.
    \item \textbf{Mô hình quản lý mẹ -- con và ma trận linh hoạt:} Mô hình này cho phép FPT kết hợp hài hòa giữa tính tập trung (kiểm soát tài chính, chuẩn hóa chính sách tại công ty mẹ) và tính phân quyền tự chủ (các công ty thành viên SBU như FPT Software, FPT Telecom chủ động kinh doanh theo thị trường ngách).
    \item \textbf{Hệ thống kiểm soát và quản trị rủi ro bài bản:} Việc áp dụng tiêu chuẩn chất lượng quốc tế nghiêm ngặt (ISO 9001, CMMi Level 5) cùng sự hiện diện của bộ phận Kiểm toán Nội bộ độc lập giúp FPT kiểm soát tốt rủi ro tuân thủ khi vận hành trên phạm vi toàn cầu.
    \item \textbf{Chiến lược đầu tư nhân lực tự chủ và bền vững:} Việc chủ động xây dựng Trường Đại học FPT và hệ sinh thái giáo dục FPT Education tạo ra nguồn cung nhân lực kỹ thuật chất lượng cao liên tục, giúp FPT giảm bớt rủi ro thiếu hụt nhân sự trên thị trường lao động.
\end{enumerate}

\subsection{Nhận xét điểm chưa đạt (Hạn chế, điểm nghẽn) cần khắc phục}
Bên cạnh những thành tựu vượt bậc, quy mô phát triển nóng cũng đặt FPT trước nhiều điểm nghẽn nghiêm trọng:

\begin{enumerate}
    \item \textbf{Bộ máy phình to, độ trễ ra quyết định gia tăng:} Với quy mô nhân sự toàn hệ sinh thái lên tới gần 90.000 người trải dài trên gần 30 quốc gia, bộ máy hành chính nội bộ xuất hiện hiện tượng nhiều tầng nấc phê duyệt trung gian, làm giảm tốc độ phản ứng trước các cơ hội kinh doanh mới nổi.
    \item \textbf{Hiện tượng ``Silo dữ liệu'' giữa các công ty thành viên:} Mặc dù thuộc cùng một tập đoàn, các công ty thành viên (FPT Telecom, FPT Retail, FPT IS, FPT Software) vẫn hoạt động như những ốc đảo dữ liệu độc lập. Sự thiếu vắng cơ chế chia sẻ tài nguyên dữ liệu người dùng làm giảm hiệu quả bán chéo sản phẩm (Cross-selling) và hạn chế sức mạnh cộng hưởng hệ sinh thái.
    \item \textbf{Áp lực cạnh tranh và biến động nhân sự công nghệ cao:} Lĩnh vực AI, dữ liệu lớn và chip bán dẫn chứng kiến sự cạnh tranh khốc liệt từ các tập đoàn công nghệ toàn cầu. Tỷ lệ luân chuyển nhân sự (turnover rate) trong ngành phần mềm ở mức cao gây áp lực nặng nề lên chi phí tuyển dụng và đào tạo lại.
    \item \textbf{Rủi ro an ninh mạng, rò rỉ dữ liệu và tuân thủ quốc tế:} Việc vận hành các trung tâm dữ liệu quy mô lớn và ứng dụng sâu rộng các mô hình AI tạo sinh (GenAI) tiềm ẩn nguy cơ lộ lọt bí mật kinh doanh của khách hàng quốc tế, vi phạm các đạo luật khắt khe như GDPR của Liên minh châu Âu.
    \item \textbf{Sự phụ thuộc lớn vào thế hệ sáng lập và bài toán kế thừa:} Uy tín cá nhân và dấu ấn cá nhân của các nhân vật thế hệ sáng lập (tiêu biểu là Chủ tịch Trương Gia Bình) vẫn mang tính chi phối sâu sắc. Bài toán quy hoạch và chuyển giao quyền lực cho thế hệ lãnh đạo kế cận (Succession Planning) là một thách thức lớn trong trung và dài hạn.
    \item \textbf{Rủi ro phụ thuộc kinh tế toàn cầu và tỷ giá:} Doanh thu xuất khẩu phần mềm phụ thuộc lớn vào thị trường Nhật Bản và Hoa Kỳ. Sự suy thoái kinh tế toàn cầu hoặc biến động trượt giá mạnh của đồng Yên Nhật (JPY) và Đô la Mỹ (USD) tác động trực tiếp đến biên lợi nhuận của tập đoàn.
\end{enumerate}

\subsection{Gợi ý và giải pháp cải thiện cho Công ty Cổ phần FPT}
Dựa trên nền tảng lý luận môn Khoa học Quản lý và Quản trị, nhóm nghiên cứu đề xuất các nhóm giải pháp nhằm nâng cao hiệu quả vận hành của FPT:

\subsubsection{1. Tinh gọn và phẳng hóa cơ cấu quản lý (Lean Management)}
\begin{itemize}
    \item Cắt giảm các cấp phê duyệt trung gian không mang lại giá trị gia tăng; trao quyền tự chủ tài chính và vận hành mạnh mẽ hơn cho các Giám đốc Chuyển giao dự án (Delivery Manager) và Trưởng nhóm kỹ thuật (Tech Lead).
    \item Ứng dụng mô hình quản lý linh hoạt (Agile Enterprise) trên quy mô toàn tập đoàn để tăng tốc độ phản hồi khách hàng.
\end{itemize}

\subsubsection{2. Xây dựng nền tảng dữ liệu dùng chung (Enterprise Data Lakehouse)}
\begin{itemize}
    \item Hiện thực hóa chiến lược \textit{DC5 (Digital Conglomerate 5.0)} bằng cách thiết lập nền tảng tích hợp dữ liệu khách hàng dùng chung (Customer 360) giữa FPT Telecom, FPT Shop, Long Châu và VnExpress.
    \item Phá vỡ các silo thông tin, tạo điều kiện cho các thuật toán AI phân tích hành vi tiêu dùng và mở ra các gói dịch vụ tích hợp đột phá.
\end{itemize}

\subsubsection{3. Hoàn thiện Khung Quản trị AI (AI Governance Framework)}
\begin{itemize}
    \item Ban hành bộ quy chuẩn nội bộ về sử dụng AI có trách nhiệm (Responsible AI), thiết lập tường lửa bảo vệ mã nguồn độc quyền của khách hàng khi kỹ sư ứng dụng LLM.
    \item Kiểm toán định kỳ các thuật toán AI để loại bỏ thiên kiến và đảm bảo tính minh bạch giải trình (Explainable AI).
\end{itemize}

\subsubsection{4. Chiến lược phát triển và giữ chân nhân tài công nghệ}
\begin{itemize}
    \item Đổi mới chính sách đãi ngộ bằng cách mở rộng chương trình phát hành cổ phiếu thưởng cho người lao động (ESOP) gắn chặt với hiệu quả OKR dài hạn.
    \item Đầu tư ngân sách đào tạo lại (Reskilling) cho toàn bộ nhân sự để chuyển đổi từ kỹ sư lập trình truyền thống sang kỹ sư AI/Cloud/Bán dẫn.
\end{itemize}

\subsubsection{5. Hoàn thiện kế hoạch chuyển giao thế hệ lãnh đạo (Succession Planning)}
\begin{itemize}
    \item Xây dựng lộ trình bổ nhiệm và chuyển giao quyền lực bài bản, minh bạch cho đội ngũ nhân sự trẻ thế hệ 8X, 9X vào các vị trí trọng yếu tại HĐQT và Ban Điều hành.
    \item Tăng cường vai trò tư vấn độc lập của các Ủy ban trực thuộc HĐQT nhằm giảm bớt sự phụ thuộc vào ý chí chủ quan của một vài cá nhân sáng lập.
\end{itemize}
